\section{The unbinned OmniFold procedure}
\label{sec:method}

\subsection{Algorithm}
\label{sec:method-algo}

OmniFold~\cite{Andreassen:2019cjw,Andreassen:2021zzk} unfolds by iterated
reweighting with machine-learned likelihood ratios, operating on per-event
observables rather than on a fixed histogram binning. Each iteration has two
steps:
\begin{enumerate}[itemsep=2pt]
  \item \textbf{Step 1 (detector space).} A classifier is trained to separate
        the measured data from the reconstructed simulation (the latter
        weighted by the current push-weights). Its calibrated output $p$
        yields the likelihood-ratio reweight $w = p/(1-p)$ on every
        reconstructed simulated event. The weight is \emph{pulled back} to
        truth level through the event pairing: each simulated event is a
        (truth, reco) pair, so the reco-space weight attaches directly to the
        truth record. Truth events with no reconstructed partner (the
        ``misses'') receive a step-1 weight from a regressor trained on the
        reconstructed subsample, $w(\mathrm{truth\ features})$ --- the native
        OmniFold acceptance treatment, with no separate efficiency
        correction.
  \item \textbf{Step 2 (truth space).} A second classifier is trained to
        separate the original truth-level simulation from the same sample
        carrying the pulled-back step-1 weights, yielding a smooth truth-space
        reweight that becomes the next iteration's push-weight.
\end{enumerate}
After $n$ iterations the push-weights define an efficiency-corrected,
truth-level reweighted simulation ensemble; histogramming it (in any binning,
chosen after the fact) gives the unfolded spectrum. The regularizer is the
iteration count together with the model capacity of the learners.

\paragraph{Background subtraction (the one binned step).}
The measured sample entering step 1 is the full set of selected data events
--- none is discarded and no negative-weight events are injected. Backgrounds
are instead removed by a per-reco-bin \emph{purity} weight, applied event by
event: the POT-scaled
background prediction (genuine backgrounds plus out-of-phase-space signal,
\S\ref{sec:selection}) is binned on the reconstructed analysis grid, and each
data event in reco bin $b$ is weighted by
\begin{equation*}
  w_{\mathrm{pur}}(b) \;=\;
  \max\!\bigl(0,\; N_{\mathrm{data}}(b) - N_{\mathrm{bkg}}(b)\bigr)
  \big/ N_{\mathrm{data}}(b),
\end{equation*}
the predicted signal purity of that bin (floored at zero in the rare bins
where the prediction exceeds the data). The step-1 classifier still sees the
full unbinned event kinematics; the analysis binning enters only through this
multiplicative correction, an ${\sim}\SI{3}{\percent}$ effect overall that is
largest where the background concentrates (\S\ref{sec:selection}).

This correction is, deliberately, the one place the analysis binning enters
before the final histogramming. Published unbinned measurements have so far
avoided the problem rather than solved it: H1 cut backgrounds to a negligible
level, LHCb corrected per track upstream of the unfold, and ATLAS left its
$<\SI{0.25}{\percent}$ background in place and assigned it as an uncertainty,
while identifying negatively-weighted simulated events as the natural
in-method subtraction~\cite{ATLAS:2024jets,Andreassen:2021zzk};
Ref.~\cite{unbinnedguide} recommends that route for nontrivial backgrounds
but records no analysis that has yet used it. The explicit per-reco-bin
weight used here is a conservative middle ground at our sub-percent
genuine-background scale. Its fully unbinned realization --- injecting the
background simulation onto the measured side with negative weight, which
targets the identical background-subtracted density $(D-B)/S$ and reproduces
the binned correction to \SI{\nwPctTot}{\percent} on the total 2D cross
section, and its systematic and statistical covariances to within
\SI{2}{\percent} (App.~\ref{app:negweight}) --- removes the residual binning dependence
and is the natural treatment for the full-phase-space extension, where
backgrounds are larger.

OmniFold has moved from method paper to production measurement over the last
few years. H1 published the first real-data applications --- jet substructure
in high-$Q^{2}$ deep-inelastic scattering, followed by a multi-differential
eight-observable jet measurement~\cite{H1:2021wkz,H1:2023fzk}; LHCb unfolded
identified-hadron distributions in $Z$-tagged jets~\cite{LHCb:2022rky}; and
ATLAS performed a simultaneous unbinned measurement of twenty-four $Z+$jets
observables~\cite{ATLAS:2024jets}. Ref.~\cite{unbinnedguide} collects these
and other applications and distills cross-experiment practice; several of the
practices audited there appear in our uncertainty budget (the ensemble-mean
cross-check of the central value, the bootstrap statistical block, and the
unified-throw covariance check of \S\ref{sec:uq-combine}).

Closest to this analysis is the study of Huang \emph{et
al.}~\cite{Huang:2024omnifold}, the first application of OmniFold to a
neutrino cross section: a mock-data unfolding of muon kinematics on the
public T2K near-detector simulated dataset, benchmarked against traditional
binned approaches. Two of its findings shape expectations for neutrino
applications. First, iteration count: Huang \emph{et al.} select 45
iterations for the OmniFold-family methods (20 for their binned benchmark)
by a binned $\chi^{2}$ against the known truth of their mock data --- a
criterion they note is not feasible on real data. We do not adopt that
recommendation, for two reasons. Their mock data injects a large
prior--data \emph{shape} discrepancy (an RPA-modified quasielastic model)
that many iterations must overcome, whereas our prior already describes the
data shapes well (\S\ref{sec:controlplots}) and the data pull is dominantly
a normalization; and cross-experiment production practice is few iterations
(typically ${\sim}5$~\cite{unbinnedguide}). Our $n=5$ is instead justified
empirically: doubling to $n=10$ moves the total cross section by
\SI{0.026}{\percent} (\S\ref{sec:method-config}), and the unfold reproduces
an independently published binned measurement of the same data
(\S\ref{sec:results}). Second, an ensemble of reweighting trials with
different train/test splits controls the ML stochastic component --- which we
budget explicitly via the seedscan band (\S\ref{sec:systematics}). The
present work differs in unfolding \emph{real} MINERvA data to an absolutely
normalized cross section with a complete uncertainty budget, in using GBDT
learners suited to a low-dimensional tabular feature set, and in validating
the method directly against a published binned measurement of the same data
(\S\ref{sec:results}).


\subsection{Estimators and configuration}
\label{sec:method-config}



We use gradient-boosted decision trees (GBDTs) rather than the neural networks
of the original OmniFold papers. The choice follows the character of the
inputs: the scalar unfolds of this note use at most five high-level features
per event, a tabular regime in which tree ensembles typically match or
outperform neural networks~\cite{Grinsztajn:2022tabular} while training
deterministically and far faster on CPU. Ref.~\cite{Milton:2025tools}
demonstrates this regime specifically for unbinned unfolding --- boosted
decision trees are effective with a small number of high-level features,
complementing deep models that ingest the full event record --- and our own
measurements follow the same division of labour: GBDTs for the scalar
unfolds, a point-cloud transformer when the input is the full cluster record
(\S\ref{sec:pet}). \S\ref{sec:validation} validates the choice directly with
calibration tests and MLP/transformer cross-checks. The engine
(\texttt{unbinned\_unfolding/python/omnifold.py}; component terminology in
App.~\ref{app:codebase}) supports interchangeable
backends; the production configuration is \textbf{LightGBM} with
$100$ trees, $8$ leaves, learning rate $0.1$ (identical settings for the
step-1 classifier, step-2 classifier, and miss regressor). Two alternative
backends --- \texttt{exact} (sklearn \texttt{GradientBoosting}) and
\texttt{hist} (sklearn \texttt{HistGradientBoosting}, matched 8-leaf
configuration) --- are retained for method-dependence studies
(\S\ref{sec:results-tension}). The production iteration count is
$n=5$; doubling to $10$ moves the total cross section by
\SI{0.026}{\percent} (\S\ref{sec:validation}). The central value is a
single-run unfold with pinned seeds; an ensemble-mean central value (averaging
independent reweighting trials) agrees with it to \SI{0.28}{\percent}, and ML
stochasticity is budgeted explicitly via the seedscan band
(\S\ref{sec:systematics}).

\paragraph{Hyperparameter choice.}
The learner settings were fixed \emph{a priori} rather than tuned on data:
100 trees, learning rate $0.1$, and depth-3 capacity are the sklearn
\texttt{GradientBoosting} defaults, and the production LightGBM backend is
configured to the matched capacity ($\texttt{num\_leaves}=8=2^{3}$) so that
every backend shares the same, deliberately modest, model capacity ---
capacity being one of OmniFold's regularisers (\S\ref{sec:method-algo}). No
hyperparameter search was performed on data, avoiding any tuning-on-result
loop. The configuration is instead validated after the fact: closure and
classifier-calibration tests (\S\ref{sec:validation}), insensitivity to
doubling the iteration count (quoted above), backend agreement in the
method-dependence studies (\S\ref{sec:results-tension}), and an explicit
seedscan band for residual ML stochasticity (\S\ref{sec:systematics}).


\subsection{MINERvA driver and pipeline conventions}
\label{sec:method-driver}



The MINERvA driver is \texttt{2d-unfolding/unfold\_2d\_omnifold\_unbinned.py};
the higher-dimensional driver
(\texttt{nd-unfolding/unfold\_nd\_omnifold\_unbinned.py}) generalizes it to an
axis list whose first two entries are always $(\pT,\pPar)$ (carrying the
fiducial gate and the per-$\pT$ flux) and \texttt{import}s the same 2D
helpers. The unfolded observables are the muon $(\pT,\pPar)$, extended by
$\Eavail$, $q_{3}$, $W$ in the higher-dimensional runs; the phase space is
that of \S\ref{sec:selection}. The following conventions hold in every
production, universe-sweep, and closure run:


\begin{itemize}[itemsep=2pt]
  \item \textbf{Reco-space masking.} The measured events entering step-1
        training are exactly those passing the reconstructed-level
        selection; nothing outside the reconstructed phase space reaches the
        classifier.
  \item \textbf{Fake-free step-1 training.} The step-1 target is the
        background-subtracted data of the purity-weight construction of
        \S\ref{sec:method-algo}: genuine backgrounds \emph{and}
        out-of-phase-space signal fakes are removed from the data side, and
        the simulated reco sample contains no fakes by construction --- so
        neither side of the step-1 training contains events whose truth lies
        outside the measurement phase space.
  \item \textbf{No double efficiency correction.} When the unfolded sample
        is histogrammed, each simulated truth event carries the product of
        its final OmniFold push-weight and its central-value simulation
        weight (the MnvTune reweighting rides along). No additional
        efficiency-histogram division is applied: the step-1 miss regression
        already returns an efficiency-corrected truth spectrum, and dividing
        again would double-correct --- a failure mode that occurred once in
        development and was caught by the marginal self-validation gate
        (\S\ref{sec:validation}).
  \item \textbf{Closure mirrors production.} In closure tests the
        pseudo-data are simulated events passing both the reconstructed and
        truth selections, weighted identically to the production data path,
        so closure exercises the same code path as the measurement rather
        than an idealised shortcut.
\end{itemize}

\subsection{Cross-section extraction}
\label{sec:method-xsec}
Unfolded truth-level counts are converted to the differential cross section
via
\begin{equation*}
  \frac{d^{n}\sigma}{d x_{1}\cdots d x_{n}}
  = \frac{U}{c\,\Phi\,N\,\mathrm{POT}\,\Delta x_{1}\cdots\Delta x_{n}}
    \times 10^{4},
\end{equation*}
where $U$ is the unfolded yield; $c$ is the per-bin \emph{completeness}, the
fraction of the truth-level sample visible to the unfolding input. Because
the event loop appends the truth-only misses explicitly, every truth event
enters the unfolding and $c\equiv1$ bin-by-bin; the factor survives in the
formula as a monitored self-check (a configuration that dropped the misses
would surface as $c<1$). $\Phi$ is the per-$\pT$ flux
integral (broadcast over the other axes), $N$ the fiducial nucleon count
($3.2353\times10^{30}$), and POT the data exposure ($1.057\times10^{21}$).
The factor $10^{4}$ converts the flux tabulation per square meter to the
reported $\si{cm^2}$ ($\SI{1}{m^2}=10^{4}\,\si{cm^2}$), so the result
carries the conventional units of $\si{cm^2}$ per axis unit per nucleon.  The implementation is
\texttt{extract\_cross\_section\_2d} and its N-dimensional generalisation
\texttt{xsec\_nd.py} (self-tested to reproduce the frozen lower-dimensional
extractions to $<10^{-12}$). The flux normalization enters as
$\sigma\propto1/\Phi$, so $\Phi$ is itself varied per flux universe when
the systematic covariance is built; this variation dominates the Flux band
(\S\ref{sec:systematics}).

\subsection{Binning}
\label{sec:method-binning}
The reported 2D binning (14 $\pT$ $\times$ 16 $\pPar$ bins; 205 of 224
reported) follows \texttt{minerva\_paper\_anc/bin\_mapping.txt}; the 19
unreported diagonal cells have
$\pT^{\mathrm{hi}}/\pPar^{\mathrm{lo}}>\tan20^{\circ}$. The extension axes use
$\Eavail$: $[0,0.1,0.2,0.4,0.8,1.5,3.0,100]\,\si{GeV}$;
$q_{3}$: $[0,0.2,0.4,0.6,0.8,1.2,2.0,100]\,\si{GeV}$;
$W$: $[0,1.1,1.4,1.8,2.2,3.0,100]\,\si{GeV}$ --- each ending in a wide catch
bin so the marginals capture the full CC-inclusive tails (a truncated top edge
would break the marginalisation anchors of \S\ref{sec:3d-anchor}). Because the
binning is applied only after unfolding, these choices can be revisited
without re-running the unfold. The full edge table and reported-bin mapping
are given in App.~\ref{app:statmethods} (\S\ref{sec:205bins}).

Neither underflow nor overflow appears in the reported product, by
construction rather than by truncation. On the truth side, the muon-axis
edges coincide with the measurement phase space itself
(\S\ref{sec:selection}) --- beyond them there is no signal definition to
measure --- and the hadronic axes span $[0,100]\,\si{GeV}$, the catch-bin
edge sitting far beyond the kinematic reach at
$\langle E_\nu\rangle\sim\SI{6}{GeV}$, so the truth grids are exhaustive. On
the reco side, a selected event whose reconstructed kinematics fall outside
the grid is treated as unaccepted: it is excluded from the measured sample
by the masking convention of \S\ref{sec:method-driver}, and its
in-phase-space truth content is recovered by the miss treatment exactly as
for any other unreconstructed event. Because the unfold is unbinned, there
is no response-matrix under/overflow to track --- the binning is applied
once, at reporting, to events that all lie inside the grids. 

Although OmniFold uses no response matrix, the detector smearing the unfold
must undo is conveniently summarized by the truth$\to$reco migration on the
analysis binning (Fig.~\ref{fig:migration}): the diagonal purity is
$\sim0.6$ (median) on every axis, i.e.\ bin-to-bin migrations are
substantial and a bin-by-bin correction would be badly biased --- the
regime requiring unfolding.

The structure differs qualitatively between the muon and hadron sides.
The spectrometer muon measurement makes the $\pT$ and $\pPar$ maps nearly
tridiagonal --- migrations rarely exceed one bin, and the purity dips
(to ${\sim}0.45$--$0.5$ for $\pT$ and ${\sim}0.35$ for $\pPar$) occur
where the binning is finest, recovering in the wide edge bins. The
calorimetric axes are coarser-binned but relatively broader and visibly
asymmetric: $\Eavail$ feeds preferentially \emph{downward} (energy carried
by neutrons and other unresponsive particles is not collected), with the
diagonal purity falling from $0.61$ in the lowest bins to $0.43$ in the
highest; $q_{3}$, built from $\Eavail$ and the muon, inherits a milder
version of the same pattern; and $W$, which compounds the $\Eavail$ and
muon smearing, is the broadest, with two-bin migrations common in the
resonance region while the lowest-$W$ (quasi-elastic) bin stays the purest
at $0.71$.

\begin{figure}[H]
  \centering
  \includegraphics[width=0.99\textwidth]{migration_resolution}
  \caption{Truth$\to$reco migration of the simulated signal on the analysis
  binning (top: column-normalized maps; bottom: diagonal purity per truth
  bin) for the five observables. Median diagonal purity is
  $0.58$--$0.61$ per axis. Shown for illustration only --- the unfold itself
  is unbinned.}
  \label{fig:migration}
\end{figure}
